% ============================================================
\chapter{M\'ethodes Proximales}
\label{ch:proximales}
% ============================================================

\section{Introduction}

Dans les ann\'ees 1960, Jean-Jacques Moreau, math\'ematicien \`a Montpellier, introduit l'\emph{op\'erateur proximal}~: pour une fonction convexe $f$, le proximal de $x$ est le point qui minimise $f$ plus un terme de p\'enalit\'e quadratique autour de $x$. L'id\'ee semble anodine, mais elle r\'ev\`ele une puissance consid\'erable~: l\`a o\`u la descente de gradient exige la diff\'erentiabilit\'e, l'op\'erateur proximal fonctionne pour toute fonction convexe semi-continue inf\'erieurement, m\^eme non lisse. Cette observation a ouvert la voie aux \emph{m\'ethodes proximales}, qui dominent aujourd'hui l'optimisation appliqu\'ee~: ISTA et FISTA pour le LASSO en statistique, ADMM pour l'optimisation distribu\'ee, l'algorithme de Douglas--Rachford pour la faisabilit\'e. Chacun exploite la d\'ecomposition d'un probl\`eme complexe en sous-probl\`emes proximaux simples.

\begin{intuition}
L'op\'erateur proximal est une g\'en\'eralisation de la projection sur un
ensemble convexe. L\`a o\`u la descente de gradient fait un pas puis projette,
l'op\'erateur proximal combine les deux en une seule \'etape~: il cherche un
point proche de l'argument qui minimise aussi la fonction. C'est un compromis
entre rester pr\`es du point courant et diminuer la fonction objectif.
\end{intuition}

\section{Op\'erateur proximal}

\begin{definition}[Op\'erateur proximal]
\index{op\'erateur proximal}
Soit $g : \R^n \to \R \cup \{+\infty\}$ convexe, s.c.i., propre.
L'\textbf{op\'erateur proximal} de $g$ est~:
\[
  \mathrm{prox}_g(v) = \argmin_{x \in \R^n} \left\{
    g(x) + \frac{1}{2} \norm{x - v}^2
  \right\}.
\]
\end{definition}

\begin{proposition}[Existence et unicit\'e]
Pour toute fonction $g$ convexe, s.c.i., propre, $\mathrm{prox}_g(v)$ existe
et est unique pour tout $v \in \R^n$.
\end{proposition}

\begin{theoreme}[Caract\'erisation]
$x^* = \mathrm{prox}_g(v)$ si et seulement si~:
\[
  v - x^* \in \partial g(x^*),
\]
o\`u $\partial g$ d\'esigne le sous-diff\'erentiel de $g$.
\end{theoreme}

\begin{proof}
La condition d'optimalit\'e de $\min_x g(x) + \frac{1}{2}\norm{x-v}^2$
s'\'ecrit~: $0 \in \partial g(x^*) + (x^* - v)$, soit
$v - x^* \in \partial g(x^*)$.
\end{proof}

\begin{formules}[title=\textbf{Op\'erateurs proximaux classiques}]
\begin{itemize}
  \item $g = 0$~: $\mathrm{prox}_g(v) = v$.
  \item $g = \iota_C$ (indicatrice d'un convexe $C$)~:
        $\mathrm{prox}_g(v) = \Pi_C(v)$ (projection).
  \item $g = \lambda \norm{\cdot}_1$ (LASSO)~:
        $[\mathrm{prox}_g(v)]_i = \mathrm{sign}(v_i)
        \max(\abs{v_i} - \lambda, 0)$ (seuillage doux).
  \item $g = \frac{\lambda}{2}\norm{\cdot}^2$~:
        $\mathrm{prox}_g(v) = \frac{v}{1 + \lambda}$.
  \item $g = \lambda \norm{\cdot}_2$ (group LASSO)~:
        $\mathrm{prox}_g(v) = v \max\!\left(1 - \frac{\lambda}{\norm{v}},
        0\right)$.
\end{itemize}
\end{formules}

\begin{definition}[R\'esolvante de Moreau]
\index{r\'esolvante de Moreau}
L'identit\'e de \textbf{Moreau} relie l'op\'erateur proximal de $g$ \`a
celui de sa conjugu\'ee $g^*$~:
\[
  \mathrm{prox}_g(v) + \mathrm{prox}_{g^*}(v) = v.
\]
\end{definition}

\section{M\'ethode de gradient proximal}

\begin{definition}[Probl\`eme composite]
On consid\`ere le probl\`eme~:
\[
  \min_{x \in \R^n} F(x) = f(x) + g(x),
\]
o\`u $f$ est convexe diff\'erentiable ($\nabla f$ est $L$-lipschitzienne)
et $g$ est convexe, s.c.i., propre (potentiellement non diff\'erentiable).
\end{definition}

\begin{algorithme}[title=\textbf{Gradient proximal (ISTA)}]
\index{ISTA}
\begin{enumerate}
  \item Choisir $x_0$, pas $\gamma \in (0, 1/L]$.
  \item Pour $k = 0, 1, 2, \ldots$~:
  \[
    x_{k+1} = \mathrm{prox}_{\gamma g}\!\left(x_k - \gamma \nabla f(x_k)\right).
  \]
\end{enumerate}
\end{algorithme}

\begin{theoreme}[Convergence d'ISTA]
Avec $\gamma = 1/L$, l'algorithme ISTA v\'erifie~:
\[
  F(x_k) - F(x^*) \leq \frac{L \norm{x_0 - x^*}^2}{2k}.
\]
La convergence est en $\mathcal{O}(1/k)$.
\end{theoreme}

\section{Acc\'el\'eration de Nesterov : FISTA}

\begin{algorithme}[title=\textbf{FISTA (Fast ISTA)}]
\index{FISTA}
\begin{enumerate}
  \item Choisir $x_0 = y_0$, $t_0 = 1$, $\gamma = 1/L$.
  \item Pour $k = 0, 1, 2, \ldots$~:
  \begin{align*}
    x_{k+1} &= \mathrm{prox}_{\gamma g}\!\left(y_k - \gamma \nabla f(y_k)\right), \\
    t_{k+1} &= \frac{1 + \sqrt{1 + 4 t_k^2}}{2}, \\
    y_{k+1} &= x_{k+1} + \frac{t_k - 1}{t_{k+1}}(x_{k+1} - x_k).
  \end{align*}
\end{enumerate}
\end{algorithme}

\begin{theoreme}[Convergence de FISTA]
L'algorithme FISTA v\'erifie~:
\[
  F(x_k) - F(x^*) \leq \frac{2L \norm{x_0 - x^*}^2}{(k+1)^2}.
\]
La convergence est en $\mathcal{O}(1/k^2)$, ce qui est \textbf{optimal}
parmi les m\'ethodes de premier ordre.
\end{theoreme}

\begin{attention}[title=\textbf{FISTA n'est pas monotone}]
Contrairement \`a ISTA, la suite $F(x_k)$ g\'en\'er\'ee par FISTA n'est pas
n\'ecessairement d\'ecroissante. L'extrapolation peut temporairement augmenter
la valeur de l'objectif.
\end{attention}

\section{ADMM}

\begin{definition}[ADMM]
\index{ADMM}
L'\textbf{Alternating Direction Method of Multipliers} r\'esout~:
\[
  \min_{x, z} f(x) + g(z) \quad \text{s.c.} \quad Ax + Bz = c.
\]
Les it\'erations sont~:
\begin{align*}
  x_{k+1} &= \argmin_x \left\{ f(x) + \frac{\rho}{2}
    \norm{Ax + Bz_k - c + u_k}^2 \right\}, \\
  z_{k+1} &= \argmin_z \left\{ g(z) + \frac{\rho}{2}
    \norm{Ax_{k+1} + Bz - c + u_k}^2 \right\}, \\
  u_{k+1} &= u_k + Ax_{k+1} + Bz_{k+1} - c,
\end{align*}
o\`u $u$ est la variable duale \'echelonn\'ee et $\rho > 0$.
\end{definition}

\begin{theoreme}[Convergence d'ADMM]
Sous des hypoth\`eses de convexit\'e de $f$ et $g$ (propres, ferm\'ees,
convexes) et la faisabilit\'e du probl\`eme, ADMM converge~:
\begin{itemize}
  \item \textbf{R\'esidu primal}~: $Ax_{k} + Bz_k - c \to 0$.
  \item \textbf{R\'esidu dual}~: $\rho A^\top B(z_k - z_{k-1}) \to 0$.
  \item \textbf{Objectif}~: $f(x_k) + g(z_k) \to p^*$.
\end{itemize}
\end{theoreme}

\begin{remarque}
ADMM est particuli\`erement efficace lorsque les sous-probl\`emes en $x$ et
$z$ admettent des solutions explicites (par exemple, moindres carr\'es +
seuillage doux pour le LASSO).
\end{remarque}

\section{S\'eparation de Douglas-Rachford}

\begin{definition}[Douglas-Rachford]
\index{Douglas-Rachford}
Pour r\'esoudre $\min_x f(x) + g(x)$, l'algorithme de Douglas-Rachford
it\`ere~:
\begin{align*}
  \hat{x}_{k} &= \mathrm{prox}_f(z_k), \\
  z_{k+1} &= z_k + \mathrm{prox}_g(2\hat{x}_k - z_k) - \hat{x}_k.
\end{align*}
\end{definition}

\begin{proposition}[Lien avec ADMM]
ADMM appliqu\'e \`a $\min f(x) + g(z)$ s.c. $x = z$ est \'equivalent \`a
Douglas-Rachford appliqu\'e \`a la formulation duale.
\end{proposition}

\section{Applications}

\begin{exemple}[LASSO -- regression parcimoneuse]
\index{LASSO}
Le probl\`eme LASSO~:
\[
  \min_x \frac{1}{2} \norm{Ax - b}^2 + \lambda \norm{x}_1
\]
est un probl\`eme composite avec $f(x) = \frac{1}{2}\norm{Ax-b}^2$ (lisse,
$L = \norm{A^\top A}$) et $g(x) = \lambda \norm{x}_1$. ISTA donne~:
\[
  x_{k+1} = S_{\lambda/L}\!\left(x_k - \frac{1}{L} A^\top(Ax_k - b)\right),
\]
o\`u $S_\tau$ est le seuillage doux composante par composante.
\end{exemple}

\begin{exemple}[D\'ebruitage par variation totale]
Pour le d\'ebruitage d'un signal $y \in \R^n$~:
\[
  \min_x \frac{1}{2}\norm{x - y}^2 + \lambda \norm{Dx}_1,
\]
o\`u $D$ est la matrice de diff\'erences finies. ADMM d\'ecompose ce
probl\`eme en un sous-probl\`eme quadratique et un seuillage doux.
\end{exemple}

\section{Exercices}

\begin{exercice}[$\star$ -- Op\'erateurs proximaux]
Calculer l'op\'erateur proximal de $g(x) = \lambda \norm{x}_1$ et de
$g(x) = \iota_{[0,+\infty)^n}(x)$ (indicatrice de l'orthant positif).
\end{exercice}

\begin{exercice}[$\star\star$ -- ISTA pour le LASSO]
Impl\'ementer ISTA pour le LASSO avec $A \in \R^{50 \times 200}$,
$x^* \in \R^{200}$ creux (10 composantes non nulles), et $b = Ax^* + \epsilon$.
Tracer la convergence et comparer avec FISTA.
\end{exercice}

\begin{exercice}[$\star\star$ -- Identit\'e de Moreau]
D\'emontrer l'identit\'e de Moreau~: $\mathrm{prox}_g(v) + \mathrm{prox}_{g^*}(v) = v$.
En d\'eduire $\mathrm{prox}_{\lambda \norm{\cdot}_\infty}$.
\end{exercice}

\begin{exercice}[$\star\star\star$ -- ADMM pour le LASSO]
D\'eriver les it\'erations ADMM pour le probl\`eme LASSO. Montrer que le
sous-probl\`eme en $x$ est un syst\`eme lin\'eaire et celui en $z$ un
seuillage doux. Impl\'ementer et comparer avec FISTA.
\end{exercice}

\begin{exercice}[$\star\star\star$ -- Convergence acc\'el\'er\'ee]
Montrer que la suite $t_k$ de FISTA v\'erifie $t_k \geq (k+1)/2$ et en
d\'eduire le taux de convergence $\mathcal{O}(1/k^2)$.
\end{exercice}
